\documentclass[12pt]{article}
\usepackage[T1]{fontenc}
\usepackage[utf8]{inputenc}
\usepackage{lmodern}
\usepackage{textcomp}
\usepackage{enumitem}
\usepackage{geometry}
\geometry{margin=1in}
\begin{document}
\begin{center}
Answer Key - Biology - Graduate Level Assessment 8 \
\small (Answers are ordered exactly as in the question paper.)
\end{center}

\section*{Multiple Choice}
\begin{enumerate}[label=\arabic*.]
\item \textbf{Answer:} D
\\ \textit{Options:} A) Gustatory receptors; B) Photoreceptors; C) Auditory receptors; D) Proprioceptors; E) Thermoreceptors
\\ \textit{Note:} Correct option: D - Proprioceptors
\item \textbf{Answer:} A
\\ \textit{Options:} A) thrombin; B) glycosaminoglycans; C) collagens; D) fibroblasts
\\ \textit{Note:} Correct option: A - thrombin
\item \textbf{Answer:} B
\\ \textit{Options:} A) Vinegar has a lower concentration of hydrogen ions than hydroxide ions; B) Vinegar has an excess of hydrogen ions with respect to hydroxide ions; C) Vinegar has a higher concentration of hydroxide ions than water; D) Vinegar has no hydroxide ions; E) Vinegar does not contain hydrogen ions
\\ \textit{Note:} Correct option: B - Vinegar has an excess of hydrogen ions with respect to hydroxide ions
\item \textbf{Answer:} A
\\ \textit{Options:} A) Protects the cell from foreign invaders; B) Acts as a storage for genetic material; C) Regulates the pH level within the cell; D) Transports oxygen throughout the cell; E) Removes nitrogenous wastes
\\ \textit{Note:} Correct option: A - Protects the cell from foreign invaders
\item \textbf{Answer:} A
\\ \textit{Options:} A) their rate of mutation should be identical.; B) they should have the exact same DNA structure.; C) they should have identical physical characteristics.; D) they should share the same number of chromosomes.; E) they live in very different habitats.
\\ \textit{Note:} Correct option: A - their rate of mutation should be identical.
\end{enumerate}

\section*{True / False}
\begin{enumerate}[label=\arabic*.]
\setcounter{enumi}{5}
\item \textbf{Answer:} False
\\ \textit{Options:} A) True; B) False
\\ \textit{Note:} The analyses show that structural characteristics of a practice are not associated with uptake of a new IT facility, but that its use may be influenced by post-graduate education in the relevant clinical condition. For this diabetes system at least, practice nurse use was critical in spreading uptake beyond initial GP enthusiasts and for sustained and rising use in subsequent years.
\item \textbf{Answer:} True
\\ \textit{Options:} A) True; B) False
\\ \textit{Note:} Our data suggest that the SCL 90-R is best viewed as an indicator of unidimensional emotional distress and somatic effects of structural brain injury.
\item \textbf{Answer:} True
\\ \textit{Options:} A) True; B) False
\\ \textit{Note:} Regarding the pedigree, we discuss different modes of inheritance. The presence of consanguineous unions in this family suggests the possibility of a common ancestor and thus a recessive autosomal mode of inheritance.
\item \textbf{Answer:} False
\\ \textit{Options:} A) True; B) False
\\ \textit{Note:} According to our study, there is a great variability when different brands and models of scanners are compared directly. Furthermore, the CT scan analysis and HU evaluation appears to gather insufficient information in order to characterize and identify the composition of renal stones.
\item \textbf{Answer:} True
\\ \textit{Options:} A) True; B) False
\\ \textit{Note:} Self-reported mechanical factors associated with chronic oro-facial pain are confounded, in part, by psychological factors and are equally common across other frequently unexplained syndromes. They may represent another feature of somatisation. Therefore the use of extensive invasive therapy such as occlusal adjustments and surgery to change mechanical factors may not be justified in many cases.
\end{enumerate}

\section*{Long Form Response}
\begin{enumerate}[label=\arabic*.]
\setcounter{enumi}{10}
\item \textbf{Answer:} Spread of disease by virulent strains of E. coli, potential military applications of virulent bacteria, introduction of cancer-causing genes into human cells. Explain why this is the correct answer and provide supporting evidence. Reference key facts or concepts that support this choice.
\\ \textit{Note:} Long-form derived from MCQ; award credit for stating the correct idea and giving supporting reasoning.
\item \textbf{Answer:} Reduction of NADP+ to NADPH. Explain why this is the correct answer and provide supporting evidence. Reference key facts or concepts that support this choice.
\\ \textit{Note:} Long-form derived from MCQ; award credit for stating the correct idea and giving supporting reasoning.
\end{enumerate}

\vfill
\noindent\textit{End of Answer Key}
\end{document}